\documentclass[12pt]{article}
\usepackage[T1]{fontenc}
\usepackage[utf8]{inputenc}
\usepackage[polish]{babel}
\usepackage{lmodern}
\usepackage[a4paper,margin=2.4cm]{geometry}
\usepackage{xcolor}
\usepackage{graphicx}
\usepackage{amsmath,amssymb}
\usepackage{booktabs}
\usepackage{array}
\usepackage{enumitem}
\usepackage{hyperref}
\usepackage{tcolorbox}
\usepackage{fancyhdr}
\definecolor{WEIIPrimary}{RGB}{0,130,150}
\definecolor{WEIIPrimaryDark}{RGB}{0,86,100}
\definecolor{WEIIPrimaryLight}{RGB}{229,246,248}
\definecolor{PRzNavy}{RGB}{31,63,142}
\definecolor{RLTextDark}{RGB}{34,34,34}
\definecolor{RLGray}{RGB}{90,90,90}
\definecolor{RLLightGray}{RGB}{245,246,248}
\definecolor{RLWarning}{RGB}{184,86,0}
\definecolor{RLExample}{RGB}{52,112,64}
\definecolor{RLCodeBg}{RGB}{248,248,248}

\usepackage{listings}
\usepackage{xcolor}

\lstdefinestyle{rlpython}{%
    language=Python,
    basicstyle=\ttfamily\scriptsize,
    keywordstyle=\color{WEIIPrimaryDark}\bfseries,
    commentstyle=\color{RLGray}\itshape,
    stringstyle=\color{RLExample},
    numbers=left,
    numberstyle=\tiny\color{RLGray},
    numbersep=6pt,
    showstringspaces=false,
    breaklines=true,
    tabsize=4,
    frame=single,
    rulecolor=\color{WEIIPrimary!35},
    backgroundcolor=\color{RLCodeBg},
    captionpos=b
}

\lstset{style=rlpython}


\hypersetup{colorlinks=true,linkcolor=WEIIPrimaryDark,urlcolor=WEIIPrimaryDark,citecolor=WEIIPrimaryDark}
\setlength{\headheight}{15pt}
\pagestyle{fancy}
\fancyhf{}
\lhead{Uczenie się ze wzmocnieniem}
\rhead{WEiI PRz}
\cfoot{\thepage}

\newtcolorbox{taskbox}[1]{colback=WEIIPrimaryLight,colframe=WEIIPrimary,title=\textbf{#1}}
\newtcolorbox{notebox}[1]{colback=RLLightGray,colframe=RLGray,title=\textbf{#1}}
\newtcolorbox{warninghandout}[1]{colback=RLWarning!8,colframe=RLWarning,title=\textbf{#1}}

\newcommand{\handouttitle}[3]{%
\begin{center}
{\Large\bfseries #1}\\[0.2cm]
{\large #2}\\[0.15cm]
{\small #3}
\end{center}
\vspace{0.4cm}
\hrule
\vspace{0.6cm}
}

\newcommand{\EnableLabNumbering}{%
  \renewcommand{\thesection}{L\arabic{section}}%
  \renewcommand{\thesubsection}{\thesection.\arabic{subsection}}%
  \renewcommand{\thesubsubsection}{\thesubsection.\arabic{subsubsection}}%
}


\title{Raport projektowy\\Uczenie się ze wzmocnieniem}
\author{Karolina Działowska\\}
\date{2026/2027}

\begin{document}
\maketitle
\tableofcontents
\newpage

\section{Informacje podstawowe}

\begin{center}
\begin{tabular}{p{0.34\textwidth}p{0.56\textwidth}}
\toprule
Element & Informacja \\
\midrule
Imię i nazwisko & Karolina Działowska\\
Numer albumu & 179478 \\
Kierunek/specjalność & Informatyka/AI \\
Tytuł projektu & Don't Touch The Spikes RL \\
Link do repozytorium & \url{https://github.com/y00nthe/Don-t-Touch-The-Spikes-rl} \\
Data oddania & 14.11.2026 \\
\bottomrule
\end{tabular}
\end{center}

\section{Temat i cel projektu}

\subsection{Opis problemu}

Projekt dotyczy prostej gry 2D wzorowanej na mobilnej grze "Don't Touch the Spikes".

Gracz steruje ptakiem na planszy o~wymiarach $400 \times 600$ pikseli.
Ptak sam leci w~poziomie, od jednej bocznej ściany do drugiej, a~po
dotknięciu ściany odbija się i~leci z~powrotem. Jednocześnie ciągle
spada w~dół, bo działa na niego grawitacja. Jedyne, co może zrobić
gracz, to skok -- podbija on ptaka w~górę.

Na suficie i~podłodze zawsze są kolce. Po każdym odbiciu kolce pojawiają
się też na ścianie, do której ptak leci, w~losowych miejscach. Każde
odbicie daje 1~punkt, a~im więcej punktów, tym więcej kolców. Na planszy
pojawiają się również monety, które można zbierać.

Zadanie polega na skakaniu w~odpowiednich momentach, tak aby ptak
przeleciał przez wolne miejsce między kolcami i~nie dotknął żadnego
kolca.

\subsection{Cel agenta}

Celem agenta jest przetrwanie jak najdłużej, czyli wykonanie jak największej
liczby odbić od ścian bez kolizji z~kolcem (na ścianie bocznej, suficie lub
podłodze), przy jednoczesnym zbieraniu monet pojawiających się na planszy.
Działanie agenta uznaje się za skuteczne, gdy potrafi grać w~nieskończoność
lub osiąga wysoki i~stabilny wynik w~każdym epizodzie.

\subsection{Uzasadnienie zastosowania uczenia się ze wzmocnieniem}

\begin{itemize}
    \item \textbf{Sekwencyjne decyzje} -- w~każdym kroku czasowym agent
    decyduje, czy wykonać skok, czy pozwolić ptakowi swobodnie opadać.
    Przebieg epizodu jest wynikiem długiej sekwencji takich decyzji.
    \item \textbf{Wpływ decyzji na przyszły stan} -- skok zmienia prędkość
    pionową ptaka, a~przez to całą trajektorię lotu aż do kontaktu ze ścianą
    docelową. Decyzja podjęta w~danej klatce decyduje o~wysokości, na której
    ptak dotrze do ściany kilkadziesiąt klatek później.
    \item \textbf{Znaczenie wyniku długoterminowego} -- zła sekwencja skoków
    nie musi przynieść natychmiastowej kary, ale może doprowadzić do stanu,
    z~którego kolizja z~kolcem na ścianie bocznej, suficie lub podłodze jest
    już nieunikniona. Istotne jest więc utrzymywanie takiej wysokości, aby
    trafić w~bezpieczną lukę między kolcami na ścianie docelowej, a~nie
    tylko unikanie kolizji w~bieżącym kroku.
    \item \textbf{Zasadność agenta uczącego się} -- problem można opisać
    jako proces decyzyjny Markowa: następny stan zależy tylko od stanu
    bieżącego i~wybranej akcji, a~skutki decyzji (nagroda lub kara)
    pojawiają się często dopiero po wielu krokach. Kolce za każdym razem
    są rozmieszczone inaczej, więc agent nie może nauczyć się na pamięć
    jednej trasy -- musi umieć reagować na dowolny układ kolców. Zamiast
    ręcznie ustalać, kiedy skakać, agent uczy się tego sam, metodą prób
    i~błędów, na podstawie nagród otrzymywanych w~trakcie gry.
\end{itemize}

\section{Formalizacja problemu RL/MDP}

\subsection{Agent i środowisko}
 
Agentem jest logika podejmująca decyzję o~skoku w~dyskretnych krokach
czasowych. Środowiskiem jest fizyka gry: grawitacja, ruch poziomy ze stałą
prędkością, odbicia od ścian bocznych, generowanie kolców i~monet oraz
wykrywanie kolizji.

W~jednym kroku epizodu:
\begin{enumerate}
    \item agent otrzymuje wektor obserwacji $o_t$,
    \item agent wybiera jedną dyskretną akcję $a_t \in \{0, 1\}$,
    \item środowisko stosuje akcję (przy skoku ustawia prędkość pionową
    ptaka), a~następnie aktualizuje jego pozycję i~prędkość,
    \item środowisko sprawdza kolizje z~kolcami, sufitem i~podłogą,
    obsługuje odbicie od ściany (zwiększenie wyniku i~zaplanowanie nowych
    kolców) oraz zebranie monety,
    \item środowisko zwraca nagrodę $r_t$, nową obserwację $o_{t+1}$
    oraz informację, czy epizod się zakończył.
\end{enumerate}

Długość jednego kroku czasowego jest stała i~wynosi $\Delta t = 1/60$~s
(jedna klatka gry, ok.\ 16{,}7~ms). Nie stosuje się powtarzania akcji
przez kilka klatek (frame skip $= 1$) -- agent podejmuje decyzję
w~każdej klatce. Na początku epizodu ptak o~wymiarach $20 \times 20$~px
znajduje się w~punkcie $(190, 290)$ z~zerową prędkością pionową,
a~ruch rozpoczyna się od razu, bez oczekiwania na pierwszy skok.

\subsection{Stan lub obserwacja}

Pełny stan środowiska obejmuje m.in.\ dokładną pozycję i~prędkość ptaka
(w~obu osiach), położenie wszystkich kolców na obu ścianach, położenie
monety, bieżący wynik, licznik opóźnienia generowania kolców oraz stan
generatora liczb losowych. Agent otrzymuje wektor obserwacji będący
zwięzłą reprezentacją części tego stanu, istotnej dla podejmowania decyzji.
Pomijane są m.in.\ kolce na ścianie, od której ptak się oddala (i~tak są
usuwane po odbiciu), wynik oraz stan generatora losowego. Przyjmuje się
idealną obserwowalność: wartości w~obserwacji są dokładne, bez szumu.

\begin{center}
\begin{tabular}{p{0.30\textwidth}p{0.60\textwidth}}
\toprule
Składowa stanu/obserwacji & Znaczenie \\
\midrule
Pozycja $Y$ agenta & Znormalizowana wysokość ptaka na planszy. \\
Odległość $X$ do ściany docelowej & Znormalizowana odległość pozioma ptaka
od ściany, w~stronę której aktualnie leci. \\
Prędkość $Y$ agenta & Prędkość pionowa ptaka (wartość ujemna -- wznoszenie,
dodatnia -- opadanie). \\
Wektor zajętości slotów (14 elementów) &
Wektor binarny opisujący ścianę docelową: 1~oznacza, że w~danym slocie
znajduje się kolec, 0~-- że slot jest bezpieczny. \\
Pozycja $X$ monety & Znormalizowana współrzędna pozioma monety;
$-1$, gdy na planszy nie ma monety. \\
Pozycja $Y$ monety & Znormalizowana współrzędna pionowa monety;
$-1$, gdy na planszy nie ma monety. \\
\bottomrule
\end{tabular}
\end{center}

\subsection{Akcje}

Zbiór akcji jest dyskretny i~dwuelementowy: $\mathcal{A} = \{0, 1\}$.
Obie akcje są dozwolone w~każdym stanie. Skok nie dodaje prędkości do
bieżącej, tylko ją nadpisuje, więc wielokrotny skok w~krótkim czasie nie
powoduje kumulowania się wznoszenia.
 
\begin{center}
\begin{tabular}{p{0.30\textwidth}p{0.60\textwidth}}
\toprule
Akcja & Opis \\
\midrule
0 -- brak akcji & Ptak swobodnie opada pod wpływem grawitacji
($g = 2000$~px/s$^2$), poruszając się poziomo w~stronę ściany docelowej. \\
1 -- skok & Prędkość pionowa ptaka zostaje natychmiast ustawiona na
$v_y = -600$~px/s (ruch w~górę); ruch poziomy nie ulega zmianie. \\
\bottomrule
\end{tabular}
\end{center}
 

\subsection{Funkcja nagrody}

\begin{equation}
    r_t =
    \begin{cases}
        -100, & \text{jeżeli w kroku } t \text{ nastąpiła kolizja z kolcem
        (ściana, sufit, podłoga)}, \\
        0{,}1 + 10 \cdot c_t, & \text{w przeciwnym przypadku},
    \end{cases}
\end{equation}
gdzie $c_t \in \{0, 1\}$ przyjmuje wartość 1, jeżeli w~kroku $t$ agent
zebrał monetę.
 
\begin{itemize}
    \item \textbf{Nagroda dodatnia} -- agent otrzymuje $+0{,}1$ za każdy
    bezpieczny krok bez kolizji oraz $+10$ za zebranie monety.
    \item \textbf{Kara} -- agent otrzymuje $-100$ za kolizję z~dowolnym
    kolcem: na ścianie bocznej, suficie lub podłodze.
    \item \textbf{Kara za krok} -- nie występuje; przeciwnie, każdy krok jest
    lekko nagradzany. Dodatnia nagroda za przeżycie zapobiega
    ,,samobójstwom'' agenta, czyli celowemu kończeniu epizodu.
    \item \textbf{Ograniczenia twarde i~miękkie} -- kara za kolizję jest
    twarda: kończy epizod. Nie stosuje się ograniczeń miękkich (np.\ kar
    za przebywanie blisko kolców).
    \item \textbf{Ryzyko niepożądanych zachowań} -- prędkość pozioma ptaka
    jest stała, więc agent nie może ,,stać w~miejscu'' i~zbierać nagrody
    za przeżycie bez wykonywania odbić; czas przetrwania jest wprost
    powiązany z~liczbą odbić, czyli wynikiem w~grze. Nagroda za monetę
    odpowiada 100~krokom przetrwania, co zachęca do zbierania monet, ale
    jest dziesięciokrotnie mniejsza od kary za śmierć, dlatego ryzykowne
    pościgi za monetą nie powinny się opłacać. Duża kara za śmierć wymusza
    unikanie przeszkód.
\end{itemize}

\subsection{Warunki zakończenia epizodu}

Epizod kończy się w~stanie porażki, gdy:
\begin{itemize}
    \item prostokąt otaczający ptaka przetnie się z~obszarem kolizji
    (hitboxem) dowolnego kolca na ścianie bocznej -- hitbox kolca jest
    pomniejszony o~5~px z~każdej strony względem jego rysunku,
    \item ptak wleci w~pas kolców na suficie lub podłodze, tzn.\ górna
    krawędź ptaka znajdzie się powyżej $y = 40$~px lub dolna krawędź
    poniżej $y = 560$~px.
\end{itemize}

Nie istnieje stan celu kończący epizod sukcesem. Epizod jest obcinany
po 10\,000 krokach (ok.\ 167~s gry); obcięcie nie jest traktowane jako
stan terminalny.

\subsection{Założenia i ograniczenia środowiska}

\begin{itemize}
    \item Fizyka gry jest uproszczona: stała grawitacja, brak oporu
    powietrza, stała prędkość pozioma, idealnie sprężyste odbicie od ścian
    bocznych (zmiana kierunku ruchu poziomego bez utraty prędkości).
    \item Zakłada się idealną obserwowalność -- agent zna dokładne pozycje
    ptaka, kolców i~monety.
    \item Pomija się opóźnienia sprzętowe i~opóźnienie reakcji -- akcja
    działa natychmiast w~kroku, w~którym została wybrana.
    \item Ściana boczna jest podzielona na sloty o~wysokości 36~px
    w~pasie między sufitem a~podłogą, co daje
    $\lfloor (600 - 2 \cdot 40) / 36 \rfloor = 14$ slotów. Liczba kolców
    wynosi $\min\left(f + \lfloor \text{wynik} / 5 \rfloor,\ 10\right)$,
    gdzie $f$ zależy od poziomu trudności -- co najmniej 4~sloty zawsze
    pozostają wolne.
    \item Po każdym odbiciu kolce z~obu ścian są usuwane, a~nowe pojawiają
    się na ścianie docelowej po 0{,}3~s. W~tym czasie wektor zajętości
    slotów ściany docelowej jest zerowy.
    \item Na planszy znajduje się co najwyżej jedna moneta. Na początku
    epizodu moneta pojawia się zawsze. Kolejna pojawia się w~losowym
    miejscu przy odbiciu, gdy wynik jest wielokrotnością 3, a~poprzednia
    została zebrana.
    \item Agent trenowany jest na poziomie trudności Medium: prędkość
    pozioma 200~px/s, początkowa liczba kolców $f = 3$.
\end{itemize}

\section{Implementacja}

\subsection{Struktura repozytorium}

Należy krótko opisać strukturę repozytorium i najważniejsze pliki projektu.

\begin{verbatim}
repozytorium/
|-- src/
|-- experiments/
|-- report/
|-- results/
`-- README.md
\end{verbatim}

\subsection{Środowisko}

Należy opisać sposób implementacji środowiska, w szczególności funkcje inicjalizacji epizodu oraz wykonania kroku środowiska.

\subsection{Agent bazowy}

Należy opisać agenta losowego, heurystycznego albo inny baseline wykorzystany do porównania z agentem uczącym się.

\subsection{Agent uczący się}

Należy opisać zastosowany algorytm RL lub DRL, jego parametry oraz sposób implementacji. W przypadku Q-tab należy opisać reprezentację tablicy wartości. W przypadku DRL należy opisać wejścia modelu, wyjścia modelu, architekturę oraz sposób uczenia.

\section{Plan eksperymentów}

Należy opisać eksperymenty zaplanowane w celu oceny działania agenta.

\begin{center}
\begin{tabular}{p{0.24\textwidth}p{0.30\textwidth}p{0.30\textwidth}}
\toprule
Eksperyment & Co jest zmieniane? & Co jest mierzone? \\
\midrule
E1 & Do uzupełnienia & Do uzupełnienia \\
E2 & Do uzupełnienia & Do uzupełnienia \\
E3 & Do uzupełnienia & Do uzupełnienia \\
\bottomrule
\end{tabular}
\end{center}

W tej sekcji należy wskazać:
\begin{itemize}
    \item liczbę epizodów uczenia,
    \item liczbę powtórzeń eksperymentu albo uzasadnienie braku powtórzeń,
    \item parametry uczenia,
    \item metryki oceny,
    \item warianty środowiska, map lub funkcji nagrody,
    \item sposób porównania z agentem bazowym.
\end{itemize}

\section{Wyniki}

W tej sekcji należy przedstawić wyniki eksperymentów. Wyniki powinny być pokazane w postaci tabel, wykresów lub czytelnych zestawień.

\begin{center}
\begin{tabular}{p{0.28\textwidth}p{0.20\textwidth}p{0.20\textwidth}p{0.20\textwidth}}
\toprule
Wariant & Średni zwrot & Skuteczność & Długość epizodu \\
\midrule
Agent bazowy & Do uzupełnienia & Do uzupełnienia & Do uzupełnienia \\
Agent RL/DRL & Do uzupełnienia & Do uzupełnienia & Do uzupełnienia \\
\bottomrule
\end{tabular}
\end{center}

\begin{figure}[h]
    \centering
    \fbox{\parbox{0.82\textwidth}{\centering Miejsce na wykres procesu uczenia}}
    \caption{Przykładowy wykres procesu uczenia agenta.}
\end{figure}

\section{Analiza i wnioski}

Należy zinterpretować wyniki. W tej części należy odpowiedzieć m.in. na pytania:
\begin{itemize}
    \item Czy agent uczący się działa lepiej niż agent bazowy?
    \item Jak zmieniało się zachowanie agenta w trakcie uczenia?
    \item Które elementy środowiska były najtrudniejsze?
    \item Czy funkcja nagrody prowadziła do oczekiwanego zachowania?
    \item Czy wystąpiły przypadki niepożądanego zachowania agenta?
    \item Jakie są najważniejsze ograniczenia rozwiązania?
\end{itemize}

\section{Instrukcja uruchomienia}

Należy opisać sposób uruchomienia projektu od czystego repozytorium. Przykładowo:

\begin{verbatim}
git clone https://github.com/...
cd nazwa-repozytorium
python -m venv .venv
.venv\Scripts\activate
pip install -r requirements.txt
python src/main.py
\end{verbatim}

Należy również wskazać, jak odtworzyć eksperymenty oraz gdzie zapisywane są wyniki.

\section{Wsparcie AI}

W tej sekcji należy obowiązkowo opisać wykorzystanie narzędzi sztucznej inteligencji. Jeżeli student nie korzystał z takich narzędzi, należy wpisać jednoznaczną informację, że nie korzystano ze wsparcia AI.

\begin{center}
\begin{tabular}{p{0.13\textwidth}p{0.19\textwidth}p{0.21\textwidth}p{0.27\textwidth}}
\toprule
Data & Narzędzie/model & Cel użycia & Transkrypt / link / zakres \\
\midrule
Do uzupełnienia & Do uzupełnienia & Do uzupełnienia & Do uzupełnienia \\
Do uzupełnienia & Do uzupełnienia & Do uzupełnienia & Do uzupełnienia \\
\bottomrule
\end{tabular}
\end{center}

Należy zapisać transkrypty rozmów z modelami AI wraz z datą oraz nazwą modelu. Jeżeli rozmowa odbywała się poza narzędziami programistycznymi, należy dodatkowo udostępnić całą rozmowę poprzez link, o ile użyte narzędzie umożliwia takie udostępnienie.

Student powinien także wyjaśnić, które elementy wygenerowane lub zasugerowane przez AI zostały samodzielnie sprawdzone, zmodyfikowane albo odrzucone.

\section{Literatura}

W tej sekcji należy wskazać wykorzystane źródła: książki, artykuły naukowe, dokumentację bibliotek, materiały wykładowe oraz inne źródła użyte podczas pracy nad projektem.

\begin{thebibliography}{9}
\bibitem{sutton_barto}
R. S. Sutton, A. G. Barto, \textit{Reinforcement Learning: An Introduction}, 2nd ed., MIT Press, 2018.
\end{thebibliography}

\end{document}
